\scp{Sound changes in all \tsc{\ili{Central} Timor} languages}{05-02}
The strongest evidence for \tsc{\ili{Central} Timor} as a distinct
subgroup comes from shared *ŋ > **g.
Proto-\ili{Central} Timor **g then devoices to \it{k} in \ili{Welaun},
\ili{Central} \ili{Mambae} everywhere, and word initially
and word finally in some varieties of South \ili{Mambae}.
Examples of *ŋ > **g are given in \trf{05-04}.

\begin{table}\tcp{\tsc{\ili{Central} Timor} *ŋ > **g}{05-04}
\begin{adjustbox}{width=\textwidth}
	\begin{threeparttable}\st{2pt}
	\begin{tabular}{lllllll}\lsptoprule
gloss	&	ear	&	name	&	soul	&	ten	&	hear	&	wind\\
PMP	&	*təli\tbr{ŋ}a	&	*\tbr{ŋ}ajan	&	*suma\tbr{ŋ}əd	&	*sa-\tbr{ŋ}a-puluq	&	*də\tbr{ŋ}əR	&	**\tbr{ŋ}əlu\su{ǁ}\\ \midrule
\ili{Welaun}	&	{\hr}li\tbr{k}a-t	&	{\hr}\tbr{k}alan-aat	&	{\hr}ma\tbr{k}ar-aat	&	{\hr}sa\tbr{k}ulu	&	{\hr}lo\tbr{k}a	&	{\hr\hr}\tbr{k}olu\\
\ili{Kemak}	&	{\hr}li\tbr{g}a-r	&	{\hr}\tbr{g}ala-r	&	{\hr}sma\tbr{g}ar-aan\su{‡}	&	{\hr}sapulu\su{Δ}	&	{\hr}re\tbr{g}a	&	\hr\cg{(agi)}\\
\ili{Tokodede}	&	{\hr}teli\tbr{g}i roa\su{†}	&	{\hr}\tbr{g}alaː	&	{\hr}ma\tbr{g}ar	&	{\hr}sa\tbr{g}ulu	&	\cg{(pliʔi)}	&	{\hr\hr}\tbr{g}elu\\
NW. \ili{Mambae}	&	{\hr}tli\tbr{g}a	&		&	{\hr}sma\tbr{g}a-n	&	{\hr}sa\tbr{g}uul	&	\cg{(pliik)}	&	{\hr\hr}\tbr{g}elo\\
C. \ili{Mambae}	&	{\hr}ti\tbr{k}a-n	&		&	{\hr}suma\tbr{k}a-n	&	{\hr}sa\tbr{k}uul	&	\cg{(fliik)}	&	{\hr\hr}\tbr{k}iul-a\\
S. \ili{Mambae}	&	{\hr}teli\tbr{g}a	&	{\hr}\tbr{k}ala	&	{\hr}samaa\tbr{k}	&	{\hr}sa\tbr{g}uul	&	\cg{(fliik)}	&	{\hr\hr}\tbr{k}eul\\
		\lspbottomrule
	\end{tabular}
		\begin{tablenotes}\tnsize
			\item[†]	{\ili{Tokodede} \it{teligi roa} ‘ear’ is from \citet{Klamer2002},
								the second element \it{roa} means ‘leaf’.}
			\item[‡]	{\ili{Kemak} \it{smagar-aan} is from \ili{Atabae} \ili{Kemak} \citep{Elliot2020}.}
			\item[Δ]	{\ili{Kemak} \it{sapulu} ‘ten’ is from *sa-ŋa-puluq with loss of
								*ŋa while other languages have loss of *ap.}
			\item[{ǁ}]	{Compare \ili{Dhao} \it{ŋəlu} ‘wind’ among others (\srf{15-04-02-02}).
								PMP *haŋin > \ili{Kemak} \it{agi} ‘wind’ has *ŋ > \it{g}.}
		\end{tablenotes}
	\end{threeparttable}
\end{adjustbox}
\end{table}

Within the wider region,
denasalisation of *ŋ is attested in most Austronesian languages
of the Bomberai Peninsula (\srf{11-05-02},
\trf{11-30}) and the \tsc{East Seram} node of \tsc{Seram-Tanimbar-Bomberai}
also have denasalisation of *ŋ in word medial position (\srf{11-02}, \trf{11-08}).

Within the Timor region,
\citet[67--68]{Edwards2021} discusses the change of **ŋ > \it{k}
in \tsc{East Rote} and \tsc{Meto}.
Most instances of \tsc{Rote-Meto} **ŋ affected by this change
are not retentions of PMP *ŋ but were reintroduced after PMP
*ŋ > Proto-Rote-\ili{Meto} **n.
\citet[67]{Edwards2021} also gives some lexical data which is
exclusively shared between \tsc{Meto}
and \tsc{\ili{Central} Timor} languages
and attributes this data,
and denasalisation of **ŋ in \tsc{East Rote}
and \tsc{Meto}, to a period of prehistoric contact between \tsc{Meto}
and \tsc{\ili{Central} Timor} languages.

While there is thus some evidence that denasalisation of **ŋ
in \tsc{Meto} and \tsc{\ili{Central} Timor} is due to contact between these languages,
it is important to note the change in \tsc{Meto} (and \tsc{East
Rote}) occurred \it{after} the break-up of Proto-Rote-\ili{Meto}.
Thus, it cannot be identified as a sound change exclusively shared
between these nodes.\footnote{
		The normal reflex of PMP *ŋ in
		Proto-Rote-\ili{Meto} is **n with 21 examples.
		There are also eight cases of PMP *ŋ = Proto-Rote-\ili{Meto} **ŋ,
		and four of PMP *ŋ > Proto-Rote-\ili{Meto} **ŋg.}

In addition to *ŋ > **g,
there are a handful of changes which provide supporting evidence for \tsc{\ili{Central} Timor}.
These changes are *R > {\0},
**r; *ma-p > **mp; *ə > \it{a} in final syllables;
and perhaps *d > **r.
Each change is discussed in turn.

\tsc{\ili{Central} Timor} languages show a split of PMP *R.
One set of lexical items has *R > {\0} in all known \tsc{\ili{Central}
Timor} reflexes, while another set of lexical items attests retention of *R in some languages.
Thus, loss of *R in a subset of lexical items can be identified
as a shared innovation at the level of \tsc{\ili{Central} Timor}.
\ili{Galolen} and the \tsc{Eastern Timor} node of \tsc{Timor-Babar}
also have *R > {\0},
though in \tsc{Eastern Timor} languages it is more regular.
Cognate sets which show *R > {\0} in all \tsc{\ili{Central} Timor}
languages are given in \trf{05-05},
and examples which show retention of *R are given in \trf{05-06}.

\begin{table}[p]\centering\tcp{\tsc{\ili{Central} Timor} *R > {\0}}{05-05}
%\begin{adjustbox}{width=\textwidth}
	\begin{threeparttable} \st{4.5pt}
	\begin{tabular}{lllllll}\lsptoprule
local gl.	&	{\hr}house	&	{\hr}bone	&	{\hr}bite	&	{\hr}new	&	{\hr}bathe	&	{\hr}day\\
PMP	&	\hr*\tbr{R}umaq	&	\hr*du\tbr{R}i	&	\hr*ka\tbr{R}at	&	\hr*baqə\tbr{R}u	&	\hr*di\tbr{R}us	&	\hr*wa\tbr{R}i\su{‡}\\
P.-C. Timor	&	**uma	&	**dui	&	**kaat	&	**bəu(n)	&	**dius	&	**wairua\\ \midrule
\ili{Welaun}	&	{\hr\hr}uma	&	{\hr\hr}lui-t	&	{\hr\hr}kaat	&	{\hr\hr}founaan	&	{\hr\hr}diis	&	{\hr\hr}wainrua\\
\ili{Kemak}	&	{\hr\hr}uma	&	{\hr\hr}rui-r	&	{\hr\hr}aat	&	{\hr\hr}heun	&	{\hr\hr}ruis	&	{\hr\hr}bairua\\
\ili{Tokodede}	&	{\hr\hr}umo\su{†}	&	{\hr\hr}rui	&	{\hr\hr}kaa	&	{\hr\hr}heuː	&	{\hr\hr}diu	&	{\hr\hr}airuu-n\\
NW. \ili{Mambae}	&	\hr\cg{(pada)}	&	{\hr\hr}rui-n	&		&	{\hr\hr}heu	&	{\hr\hr}riu	&	{\hr\hr}airua-n\\
C. \ili{Mambae}	&	\hr\cg{(fada)}	&	{\hr\hr}rui-n	&		&	{\hr\hr}heun	&	{\hr\hr}lariu	&	{\hr\hr}airua-n\\
S. \ili{Mambae}	&	{\hr\hr}uum	&	{\hr\hr}rui	&	\hr\cg{(hoi)}	&	{\hr\hr}heu	&	{\hr\hr}rio	&	{\hr\hr}airua\\
		\lspbottomrule
	\end{tabular}
		\begin{tablenotes}\tnsize
			\item[†]	{\ili{Tokodede} \it{umo} is from the \ili{Maubara} variety.
								Other varieties of \ili{Tokodede} have \it{soa} ‘house’.}
			\item[‡]	{Reflexes of *waRi are compounds of *waRi + *duha ‘two’.
								\ili{Welaun} \it{wainrua} and \ili{Kemak} \it{bairua} are given as ‘day after tomorrow’.
								\ili{Tokodede} \it{airuu-n}, and \tsc{Mambae} \it{airua-n, airua} are
								given as ‘day before yesterday’.
								\ili{Tokodede} also has \it{aer-ruu} ‘day after tomorrow’
								and \it{aer} ‘tomorrow’ which may also be cognate with *waRi
								with *R = \it{r}
								and final CV > VC metathesis.
								Another cognate set related to this etymon is *pa-waRi > \ili{Welaun}
								\it{a-wai}, \ili{Kemak} \it{pai}, \ili{Tokodede} \it{pai} ‘dry in sun’.}
		\end{tablenotes}
	\end{threeparttable}
%\end{adjustbox}
\end{table}

\begin{table}[p]\centering\tcp{\tsc{\ili{Central} Timor} *R = **r}{05-06}
\begin{adjustbox}{width=\textwidth}
	\begin{threeparttable}\st{2pt}
	\begin{tabular}{llllllll}\lsptoprule
local gl.	&	{\hr}root	&	{\hr}brother	&	{\hr}red	&	{\hr}blood	&	{\hr}run	&	{\hr}water	&	{\hr}bamboo\\
PMP	&	\hr*\tbr{R}amut	&	\hr*ña\tbr{R}a	&	\hr*ma-i\tbr{R}aq	&	\hr*da\tbr{R}aq	&	\hr*la\tbr{R}iw	&	\hr*wahiyə\tbr{R}	&	\hr*qau\tbr{R}\\
P.-C. Timor	&	**\tbr{r}amut	&	**na\tbr{r}a	&	**me\tbr{r}a	&	**da\tbr{r}a	&	**pala\tbr{r}i	&	**wee\tbr{r}	&	**au\tbr{r}\\ \midrule
\ili{Welaun}	&	\hr\cg{(waʔa-n)}	&	\hr\cg{(naʔi)}	&	{\hr\hr}meak	&	{\hr\hr}laa-t	&	{\hr\hr}a-lai	&	{\hr\hr}wee	&	{\hr\hr}au\\
\ili{Kemak}	&	\hr\cg{(haʔa-n)}	&	{\hr\hr}naa-r	&	{\hr\hr}meak	&	{\hr\hr}raa-r	&	{\hr\hr}pə̆lai	&	{\hr\hr}bia, bea\su{Δ}	&	{\hr\hr}oa\\
\ili{Tokodede}	&	{\hr\hr}\tbr{r}amut	&	\hr\cg{(ambou)}	&	{\hr\hr}meoː	&	{\hr\hr}raa	&	{\hr\hr}plai	&	{\hr\hr}ee	&	{\hr\hr}betul ae\su{ǁ}\\
NW. \ili{Mambae}	&	{\hr\hr}\tbr{r}amu-n	&	{\hr\hr}na\tbr{r}a	&	{\hr\hr}me\tbr{r}a	&	{\hr\hr}la\tbr{r}a\su{†}	&		&	{\hr\hr}ee\tbr{r}-a	&	\hr\cg{(betun)}\\
C. \ili{Mambae}	&	{\hr\hr}\tbr{r}amu-n	&	{\hr\hr}na\tbr{r}a-n	&	{\hr\hr}me\tbr{r}a-n	&	{\hr\hr}la\tbr{l}a	&		&	{\hr\hr}ee\tbr{r}-a	&	\hr\cg{(betun)}\\
S. \ili{Mambae}	&	\hr\cg{(haa)}	&	{\hr\hr}na\tbr{r}a	&	{\hr\hr}me\tbr{r}a	&	{\hr\hr}la\tbr{r}a	&	{\hr\hr}flae\tbr{r}\su{‡}	&	{\hr\hr}ee\tbr{r}	&	{\hr\hr}oo\tbr{r}\\
		\lspbottomrule
	\end{tabular}
		\begin{tablenotes}\tnsize
			\item[†]	{\tsc{Mambae} has dissimilation of the initial liquid in this
								form; *daRaq > **dara > **rara > \it{lara}.
								\ili{Central} \ili{Mambae} from Liquidoe has subsequent assimilation of the
								medial liquid to yield \it{lala}.
								Other varieties of \ili{Central} \ili{Mambae} in \citet{Fogaca2017} have
								\it{lara} or \it{lara-n} ‘blood’.}
			\item[‡]	{South \ili{Mambae} \it{flaer} ‘run’ is a metathesised form.
								Because South \ili{Mambae} has lowering of /i/ after metathesis when
								the previous vowel is /a/ (e.g.
								*kami \it{ami → aem} ‘we (excl.)’),
								the unmetathesised form of \it{flaer} ‘run’ could be either **flari or **flare.}
			\item[Δ]	{Most varieties of \ili{Kemak} have \it{bia} ‘water’.
								\ili{Marobo} and \ili{Leosibe} \ili{Kemak} have \it{bia {\tl} bea}.}
			\item[{ǁ}]	{\ili{Tokodede} \it{betul ae} is glossed as \ili{Tetun} Dili \it{au
								maus}, \ili{English} ‘the small bamboo’.
								The term \it{betul} (from PMP *bətuŋ) is given as the generic
								word for bamboo \citep[4]{LSG2006}.}
		\end{tablenotes}
	\end{threeparttable}
\end{adjustbox}
\end{table}

Among the words in \trf{05-06} with retention of *R,
the sound changes are simplest in \tsc{Mambae} and \ili{Welaun}.
\tsc{Mambae} has *R = \it{r} in all positions,
while \ili{Welaun} has *R > {\0} in all positions.
The usual pattern in \ili{Tokodede} is *R > {\0},
with the exception of *Ramut > \it{ramut} ‘root’.
Additionally, *qauR > \it{betul ae} `bamboo' has fronting
of the final vowel along with loss of *R.
This may indicate that *R > {\0} in \ili{Tokodede} went through an
intermediate glide stage **y.

\ili{Kemak} has *R > {\0} word medially,
but *R > \it{a} word finally.
Evidence that the final vowel of \ili{Kemak} \it{bia,
bea} ‘water’ and \it{oa} ‘bamboo’ is indeed a reflex of *R comes
from words with a final double vowel such as *buaq > \it{boo,
buu} ‘betel nut’, *kahu > \it{oo} ‘you (sg.)’,
and *həsi > \it{sii} ‘meat,
flesh’ which do not show final \it{a}.
Another possible example of *R > \it{a} word finally in \ili{Kemak}
is *niuR > \it{nua,
noa} (\ili{Kutubaba} \ili{Kemak} \it{noa} ‘coconut’,
\ili{Marobo} \ili{Kemak} \it{noa {\tl} nua} ‘coconut’),
though in this case \tsc{Mambae} has *R > {\0},
as shown by Northwest/\ili{Central} \ili{Mambae} \it{noo-a},
South \ili{Mambae} \it{noo} ‘coconut’.

Another sound change shared between all \tsc{\ili{Central} Timor} languages
is *ə > \it{a} in most final syllables.
The main exception is when the final consonant is *q,
in which case there is sporadic raising of *ə > \it{e, i}.
Both these sound changes are also found in some \tsc{Timor-Babar}
languages spoken on the Timor mainland.
Examples of *ə in final syllables are given in \trf{05-07}. (See
also \trf{05-10} on \prf{tab:05-10} for another example of *ə > \it{e,
i} /{\gap}q{\#} in reflexes of *daRəq ‘earth’).


\begin{table}\tcp{\tsc{\ili{Central} Timor} *ə /{\gap}(C){\#}}{05-07}
%\begin{adjustbox}{width=\textwidth}
	\begin{threeparttable}\st{3pt}
	\begin{tabular}{lllllll}\lsptoprule
local gl.	&	teeth	&	black	&	inside	&	fly \cg{(n.)}	&	wash	&	foam\\
PMP	&	*nip\tbr{ə}n	&	*ma-qit\tbr{ə}m	&	*dal\tbr{ə}m\su{†}	&	*lal\tbr{ə}j	&	*bas\tbr{ə}q\su{Δ}	&	*buj\tbr{ə}q\\ \midrule
\ili{Welaun}	&	\cg{(kiis-aat)}	&	{\hr}met\tbr{a}n	&	{\hr}lam\tbr{a}raan	&	{\hr}lal\tbr{a}n	&	{\hr}sah\tbr{i}	&	{\hr}for\tbr{e}n\\
\ili{Kemak}	&	{\hr}nip\tbr{a}-r	&	{\hr}met\tbr{a}m	&	{\hr}lar\tbr{a}-n	&	{\hr}lal\tbr{a}	&	{\hr}has\tbr{a}	&	{\hr}bia hul\tbr{a}n\\
\ili{Tokodede}	&	{\hr}nip\tbr{a}	&	{\hr}met\tbr{a}:	&	{\hr}kai lar\tbr{a}\su{‡}	&	{\hr}lal\tbr{a}	&	{\hr}kahas\tbr{a}	&	{\hr}hur\tbr{i}ː\\
NW. \ili{Mambae}	&	{\hr}nip\tbr{a}-n	&	{\hr}met\tbr{a}	&	{\hr}ai lal\tbr{a}-n\su{‡}	&		&	{\hr}ha\tbr{a}s	&	\\
C. \ili{Mambae}	&	{\hr}nif\tbr{a}-n	&	{\hr}met\tbr{a}	&	{\hr}ai lal\tbr{a}-n\su{‡}	&		&	{\hr}ha\tbr{e}s	&	\\
S. \ili{Mambae}	&	{\hr}nif\tbr{a}	&	{\hr}met\tbr{a}	&	{\hr}lal\tbr{a}	&		&	{\hr}ha\tbr{a}s	&	\\
		\lspbottomrule
	\end{tabular}
		\begin{tablenotes}\tnsize
			\item[†]	{The \ili{Kemak},
								\ili{Tokodede}, and \ili{Welaun} reflexes all have metathesis of the initial
								and medial consonants: *daləm > **ralam > **laram.
								(Such metathesis may have also taken place in \tsc{Mambae},
								though assimilation of **r > \it{l} means that we cannot be sure.)
								\ili{Welaun} then has metathesis of the medial
								and final consonant to yield **laram > **lamar.
								This consonant metathesis in \ili{Welaun} must have been complete before
								the change of non-final **r > \it{l}.
								Final \it{aan} in \ili{Welaun} is a historic genitive suffix.}
			\item[‡]	{\ili{Tokodede} \it{kai lara}
								and Northwest/\ili{Central} \ili{Mambae} \it{ai lala-n} mean ‘forest’ from ‘tree’ + ‘inside’.
								\ili{Kemak} also has \it{ai lara-n} ‘forest’.}
			\item[Δ]	{\ili{Welaun} \it{sahi} ‘wash’ has consonant metathesis.
								\ili{Tokodede} \it{kahasa} ‘wash’ is from \citet{Klamer2002}.
								The unmetathesised form of \ili{Central} \ili{Mambae} \it{haes} ‘wash’ may
								be *[hasi] or *[hase].
								The unmetathesised form of South \ili{Mambae} \it{haas} ‘wash’ would probably be *[hasa].}
		\end{tablenotes}
	\end{threeparttable}
%\end{adjustbox}
\end{table}

Another change which is probably shared between all \tsc{\ili{Central}
Timor} languages involves the reflexes of *ma-p.
\ili{Welaun} has *ma-p > \it{f} while other \tsc{\ili{Central} Timor} languages
usually have *ma-p > \it{b},
though there are also sporadic cases of *ma-p > \it{p}.
(Instances of *ma-p > \it{p} may simply be reflexes of *p without the prefix *ma-).
Because these reflexes must be kept distinct from the reflexes
of *p (see \trf{05-01}, \prf{tab:05-01}),
they probably point to an intermediate nasal-stop sequence **mp
at the level of Proto-\ili{Central} Timor.
Examples of PMP *ma-p > **mp in \tsc{\ili{Central} Timor} are given
in \trf{05-08}.

\begin{table}[p]\centering\tcp{\tsc{\ili{Central} Timor} *ma-p > **mp}{05-08}
%\begin{adjustbox}{width=\textwidth}
	\begin{threeparttable}\st{2pt}
	\begin{tabular}{llllll}\lsptoprule
local gl.	&	{\hr}hot	&	{\hr}white	&	{\hr}full	&	{\hr}bitter	&	{\hr}sick\\
PMP	&	\hr*\tbr{ma-p}anas	&	\hr*\tbr{ma-p}utiq	&	\hr*\tbr{ma-p}ənuq	&	**(\tbr{ma-})\tbr{p}əju\su{ǁ}	&	**\tbr{ma-p}əjəs\su{¶}\\
P.-C. Timor	&	**\tbr{mp}anas	&	**\tbr{mp}uti	&	**\tbr{mp}ənu	&	**\tbr{mp}əju	&	**\tbr{mp}əjas\\ \midrule
\ili{Welaun}	&	{\hr\hr}\tbr{f}anas	&	\hr\cg{(baras)}	&	{\hr\hr}\tbr{h}onu\su{Δ}	&	{\hr\hr}\tbr{f}oru	&	{\hr}\tbr{f}oras\\
\ili{Kemak}	&	{\hr\hr}\tbr{b}ansan\su{‡}	&	{\hr\hr}\tbr{b}uti	&	{\hr\hr}\tbr{b}enu	&	{\hr\hr}\tbr{p}elun	&	\cg{(bə̆ragan)}\\
\ili{Tokodede}	&	{\hr\hr}\tbr{b}anaː	&	{\hr\hr}\tbr{b}uti	&	{\hr\hr}\tbr{b}enu	&	{\hr\hr}\tbr{p}elu	&	\cg{(groʔa)}\\
NW. \ili{Mambae}	&	{\hr\hr}\tbr{b}ana	&	{\hr\hr}\tbr{b}uti	&	{\hr\hr}\tbr{b}enu	&		&	\cg{(baan)}\\
C. \ili{Mambae}\su{†}	&	{\hr\hr}\tbr{p}aan	&	{\hr\hr}\tbr{b}uti	&	{\hr\hr}\tbr{p}eun	&		&	\cg{(paan)}\\
S. \ili{Mambae}	&	\hr\cg{(brusi)}	&	{\hr\hr}\tbr{b}uti	&	{\hr\hr}\tbr{b}eun	&		&	\cg{(moras)}\\
		\lspbottomrule
	\end{tabular}
		\begin{tablenotes}\tnsize
			\item[†]	{\ili{Central} \ili{Mambae} words here are from Laulara.
								The word \it{paan} is given as ‘sick’.}
			\item[‡]	{\ili{Kemak} \it{bansan} ‘hot’ is from intermediate **banas-aan with
								deletion of the second vowel.
								\ili{Marobo} \ili{Kemak} is known to have both \it{bansan} ‘hot’ and \it{banas-aan} ‘heat’.}
			\item[Δ]	{\ili{Welaun} \it{honu} ‘full’ is from *pənuq without initial *ma-.}
			\item[{ǁ}]	{**ma-pəju ‘bitter’ is from PMP *qapəju ‘gall, gallbladder, bile’.}
			\item[¶]	{**ma-pəjəs is from PMP *hapəjəs ‘smarting, stinging pain’.
								South \ili{Mambae} \it{moras} ‘sick’ is a borrowing from \ili{Tetun},
								as evidenced by irregular *ə > \it{o} rather than expected *ə > \it{e}.}
		\end{tablenotes}
	\end{threeparttable}
%\end{adjustbox}
\end{table}

The reflexes of putative **mp also have to be kept distinct from reflexes of **mb.
While both nasal-stop sequences merge as \it{b} in most \tsc{\ili{Central}
Timor} languages, in \ili{Welaun} the reflexes are different with **mb > \it{b} as opposed
to **mp > \it{f.} Examples of **mb are given in \trf{05-09}.

\begin{table}[p]\centering\tcp{\tsc{\ili{Central} Timor} **mb}{05-09}
%\begin{adjustbox}{width=\textwidth}
	\begin{threeparttable} \st{2pt}
	\begin{tabular}{ll@{}l@{}ll@{}l}\lsptoprule
local gl.	&	{\hr}owner, master	&	{\hr}carrying stick	&	{\hr}betel nut	&	saliva	&	slap\\
PMP	&	\hr*umpu	&	\hr*lemba	&	\hr*buaq	&	**kambeR\su{ǁ}	&	**mbasaR\su{¶}\\
P.-C. Timor	&	**u\tbr{mb}u	&	**kai le\tbr{mb}a\su{‡}	&	**\tbr{mb}ua\su{Δ}	&	**ka\tbr{mb}a	&	**\tbr{mb}asa\\ \midrule
\ili{Welaun}	&	{\hr\hr}u\tbr{b}un	&	{\hr\hr}ai kale\tbr{b}a	&	{\hr\hr}\tbr{b}ua	&	{\hr\hr}a\tbr{b}a-t	&	{\hr\hr}\tbr{b}asa\\
\ili{Kemak}	&	{\hr\hr}u\tbr{b}un\su{†}	&	{\hr\hr}ai le\tbr{b}an	&	{\hr\hr}\tbr{b}oo	&	{\hr\hr}a\tbr{b}a-r bea-n\su{†}	&	{\hr\hr}\tbr{b}asa\\
\ili{Tokodede}	&	{\hr\hr}u\tbr{b}u	&	{\hr\hr}kai le\tbr{b}a	&	{\hr\hr}\tbr{b}uo	&	{\hr\hr}\tbr{b}aa, ka\tbr{b}ee	&	{\hr\hr}k\tbr{b}aas\\
S. \ili{Mambae}	&	{\hr\hr}u\tbr{b}u	&		&	{\hr\hr}\tbr{b}uu	&	{\hr\hr}a\tbr{b}a	&	{\hr\hr}\tbr{b}aas\\
		\lspbottomrule
	\end{tabular}
		\begin{tablenotes}\tnsize
			\item[†]	{\ili{Kemak} \it{ubun} ‘owner,
								master’ and \it{aba-r bea-n} ‘saliva’ are from \ili{Atabae} \ili{Kemak} \citep{Elliot2020}.}
			\item[‡]	{PCMP *lemba is ‘carry on a shoulder pole’.
								Reflexes here are compounded with reflexes of *kayu ‘wood’
								and refer to the pole with which something is carried over the shoulder.}
			\item[Δ]	{PMP *buaq often develops into a form with an initial nasal-stop
								cluster in languages of Timor when the meaning is ‘betel nut’.
								Compare, for instance, Proto-Rote-\ili{Meto} *mbuah \citep[270]{Edwards2021}.}
			\item[{ǁ}]	{For **kambeR ‘saliva’ compare Proto-Rote-\ili{Meto} *kambe,
								\ili{Kisar} \it{apar,} \ili{Kamarian} \it{aper} ‘saliva’ among other cognates \citep[192]{Edwards2021}.
								\ili{Tokodede} \it{baa} is from \citet{LSG2006},
								\it{kabee} is from \citet{Klamer2002}.
								\it{kabee} may a compound of **kaba
								and \it{ee} ‘water’ or a borrowing from \ili{Tetun} \it{kabeen} ‘saliva’.}
			\item[¶]	{For **mbasaR ‘slap’ compare Proto-Rote-\ili{Meto} *mbasa,
								\ili{Helong} \it{papas}, Ili{\Q}uun \it{paas},
								\ili{Kisar} \it{pahar}, and \ili{Sika} \it{pahak} among other cognates \citep[263]{Edwards2021}.}
		\end{tablenotes}
	\end{threeparttable}
%\end{adjustbox}
\end{table}

Finally, all \tsc{\ili{Central} Timor} languages also have *d > **r.
While this results in a merger with those words which retain
*R in \tsc{Mambae}, other languages do not merge *d and *R.
Because of this, we must either assign a value other than [r]
to the Proto-\ili{Central} Timor reflex of *R (when retained),
and/or we must posit that *d > **r occurred after the break-up
of the group, and would thus provide no evidence for \tsc{\ili{Central} Timor}.
Whatever hypothesis is taken,
the evidence that *d > **r provides for \tsc{\ili{Central} Timor} would
be negligible since it is an extremely common change.

\ili{Welaun} has subsequent *d > **r > \it{l} word initially
and medially, but retains *d > **r = \it{r} word finally,
as shown in \trf{05-10}.
\tsc{Mambae} has occasional instances of *d > **r > \it{l} which
can usually be attributed to assimilation or dissimilation with another adjacent liquid.
Two examples of \tsc{Mambae} *d > \it{l} are *daRaq > **rara
> \it{lara, lala} ‘blood’ (\trf{05-06}, \prf{tab:05-06}),
and *daləm > **rala > \it{lala} ‘inside’ (\trf{05-07}, \prf{tab:05-07}).

\begin{table}\tcp{\tsc{\ili{Central} Timor} *d}{05-10}
\begin{adjustbox}{width=\textwidth}
	\begin{threeparttable}\st{2pt}
	\begin{tabular}{llllllll}\lsptoprule
local gl.	&	bone	&	earth	&	leaf	&	stand	&	young	&	knee	&	soul\\
PMP	&	*\tbr{d}uRi	&	*\tbr{d}aRəq	&	*\tbr{d}ahun	&	*\tbr{d}iRi	&	*ŋu\tbr{d}a\su{‡}	&	*tuhu\tbr{d}\su{ǁ}	&	*sumaŋə\tbr{d}\\ \midrule
\ili{Welaun}	&	{\hr}\tbr{l}ui-t	&	{\hr}\tbr{l}ai	&	{\hr}\tbr{l}oon	&	\cg{(a-fii)}	&	{\hr}ku\tbr{l}aʔaan	&	{\hr}tuu\tbr{r}-aat	&	{\hr}maka\tbr{r}-aat\\
\ili{Kemak}	&	{\hr}\tbr{r}ui-r	&	{\hr}\tbr{r}ae	&	\cg{(ai taha-n)}	&	\cg{(ara)}	&		&	{\hr}tu\tbr{r}u-r	&	{\hr}smaga\tbr{r}-aan\\
\ili{Tokodede}	&	{\hr}\tbr{r}ui	&	{\hr}\tbr{r}ae	&	{\hr}\tbr{r}oaː	&	{\hr}b\tbr{r}ii\su{Δ}	&		&	\cg{(tunuruː)}	&	{\hr}maga\tbr{r}\\
NW. \ili{Mambae}	&	{\hr}\tbr{r}ui-n	&	{\hr}\tbr{r}ai-a	&	\cg{(ai nora)}\su{†}	&		&	{\hr}gu\tbr{r}a	&		&	{\hr}smaga-n\\
C. \ili{Mambae}	&	{\hr}\tbr{r}ui-n	&	{\hr}\tbr{r}ai-a	&	\cg{(nora-n)}\su{†}	&		&	{\hr}ku\tbr{r}a	&		&	{\hr}sumaka-n\\
S. \ili{Mambae}	&	{\hr}\tbr{r}ui	&	{\hr}\tbr{r}ae	&	\cg{(ai nora)}	&	{\hr}m\tbr{r}ii	&	{\hr}kuu\tbr{r}	&	\cg{(oe tunur)}	&	{\hr}samaak\\
		\lspbottomrule
	\end{tabular}
		\begin{tablenotes}\tnsize
			\item[†]	{Northwest \ili{Mambae} has \it{ulu laun} ‘head hair’
								and \ili{Central} \ili{Mambae} has \it{ulu lahon} ‘head hair’ the second
								elements of which may be irregularly connected with *dahun ‘leaf’ (see \srf{02-03-01-06}).}
			\item[‡]	{The \tsc{Mambae} reflexes mean ‘grass’.
								The reconstructed PMP gloss is ‘young, of plants’.}
			\item[Δ]	{\ili{Tokodede} \it{brii} ‘stand’ is from \citet{Klamer2002}.}
			\item[{ǁ}]	{\ili{Tokodede} \it{tunuruː}
								and South \ili{Mambae} \it{oe tunur} ‘knee’ may be irregularly connected
								with PMP *tuhud, though the source of medial \it{n} is unknown.
								(South \ili{Mambae} \it{oe tunur} is from \citealt{Klamer2002}).}
		\end{tablenotes}
	\end{threeparttable}
\end{adjustbox}
\end{table}

To summarise, the change of PMP *ŋ > **g can be identified as
a shared innovation which defines the \tsc{\ili{Central} Timor} subgroup.
While \tsc{Meto} and \tsc{East Rote} have similar **ŋ > \it{k},
this change occurred after the break-up of the \tsc{Rote-Meto}
group and does not affect most instances of PMP *ŋ in \tsc{Rote-Meto}.
Other sound changes shared between \tsc{\ili{Central} Timor} languages
which provide supporting evidence for this group are the split
of *R > {\0},
**r in the same words,
*ə > \it{a} in word-final syllables,
*ma-p > **mp, and perhaps *d > **r.
